% DERV-04: Softmax vs ILR Compositional Closure for CF-LIBS
% ASSERT_CONVENTION: natural_units=none, unit_system=SI_eV_cm3, composition=number_fractions_sum1
%
% References:
%   Egozcue et al. (2003) Math. Geology 35(3), 279-300
%   Aitchison (1986) The Statistical Analysis of Compositional Data
%
% Codebase references:
%   cflibs/inversion/joint_optimizer.py (softmax parameterization)
%   cflibs/inversion/closure.py (ILR transform, Helmert basis)

\documentclass[11pt]{article}
\usepackage{amsmath,amssymb,amsthm}
\usepackage{booktabs}
\usepackage{hyperref}

\newcommand{\CC}{\mathbf{C}}
\newcommand{\ttheta}{\boldsymbol{\theta}}
\newcommand{\zz}{\mathbf{z}}
\newcommand{\JJ}{\mathbf{J}}
\newcommand{\VV}{\mathbf{V}}
\newcommand{\ones}{\mathbf{1}}

\title{Softmax vs.\ Isometric Log-Ratio Parameterization\\
for CF-LIBS Compositional Closure}
\author{CF-LIBS GPU Optimization Project --- DERV-04}
\date{}

\begin{document}
\maketitle

%% ====================================================================
\section{Physics Motivation}
\label{sec:motivation}
%% ====================================================================

In Calibration-Free LIBS, element concentrations are determined from
Boltzmann plot intercepts:
%
\begin{equation}
  \label{eq:raw-concentration}
  C_s \propto U_s(T)\,\exp(q_s),
\end{equation}
%
where $U_s(T)$ is the partition function and $q_s$ is the intercept
from the Boltzmann plot of element $s$.

The \emph{closure constraint} $\sum_{s=1}^D C_s = 1$ (number fractions)
eliminates the unknown experimental factor $F$ and confines the
composition to the $(D{-}1)$-dimensional probability simplex:
%
\begin{equation}
  \label{eq:simplex}
  \Delta^{D-1}
  = \bigl\{\CC \in \mathbb{R}^D : C_i \ge 0,\;\textstyle\sum_{i=1}^D C_i = 1\bigr\}.
\end{equation}
%
% DIMENSION CHECK: C_i are dimensionless number fractions. PASSED.

For gradient-based joint optimization (jointly fitting $T$, $\ne$, and
$\{C_s\}$), we need a differentiable map from unconstrained parameters
to $\Delta^{D-1}$.  Two standard approaches are:

\begin{itemize}
  \item \textbf{Softmax} (overparameterized): $\mathbb{R}^D \to \Delta^{D-1}$.
    Already implemented in \texttt{joint\_optimizer.py}.
  \item \textbf{ILR} (minimal parameterization):
    $\mathbb{R}^{D-1} \to \Delta^{D-1}$.
    Already implemented in \texttt{closure.py} via Helmert basis.
\end{itemize}


%% ====================================================================
\section{Softmax Parameterization}
\label{sec:softmax}
%% ====================================================================

\subsection{Definition}

Given unconstrained parameters $\ttheta = (\theta_1, \ldots, \theta_D) \in \mathbb{R}^D$:
%
\begin{equation}
  \label{eq:softmax}
  C_i = \frac{\exp(\theta_i)}{\sum_{j=1}^D \exp(\theta_j)},
  \qquad i = 1, \ldots, D.
\end{equation}
%
% DIMENSION CHECK: [theta_i] = dimensionless.  exp(dim-less) = dim-less.
% [C_i] = dim-less / dim-less = dimensionless.  PASSED.

\paragraph{Log-sum-exp stabilization.}
To prevent overflow for large $|\theta_i|$, define
$\theta_{\max} = \max_j \theta_j$ and compute:
%
\begin{equation}
  \label{eq:logsumexp}
  C_i = \frac{\exp(\theta_i - \theta_{\max})}
             {\sum_{j=1}^D \exp(\theta_j - \theta_{\max})}.
\end{equation}
%
This is numerically stable for $|\theta_i| < 700$ (float64 exp limit).
The shift does not change the result since
$\exp(\theta_i - c)/\sum_j \exp(\theta_j - c) = \exp(\theta_i)/\sum_j \exp(\theta_j)$
for any constant $c$.

\paragraph{Shift symmetry.}
The softmax is invariant under $\ttheta \to \ttheta + c\,\ones$:
%
\begin{equation}
  C_i(\ttheta + c\,\ones) = C_i(\ttheta) \quad \forall\,c \in \mathbb{R}.
\end{equation}
%
This means the map $\mathbb{R}^D \to \Delta^{D-1}$ is $D$-to-$(D{-}1)$:
one degree of freedom is redundant.

\subsection{Jacobian}

The Jacobian $\JJ_\mathrm{sm} \in \mathbb{R}^{D \times D}$ with entries
$J_{ij} = \partial C_i / \partial \theta_j$ is:
%
\begin{equation}
  \label{eq:softmax-jacobian}
  \boxed{
    J_{ij}^\mathrm{sm} = C_i\,(\delta_{ij} - C_j),
  }
\end{equation}
%
or in matrix form:
%
\begin{equation}
  \label{eq:softmax-jacobian-matrix}
  \JJ_\mathrm{sm} = \mathrm{diag}(\CC) - \CC\,\CC^T.
\end{equation}
%
% DIMENSION CHECK: [J_ij] = [C_i] * [delta_ij - C_j] = dim-less. PASSED.

\paragraph{Derivation.}
Let $Z = \sum_j \exp(\theta_j)$.  Then $C_i = \exp(\theta_i)/Z$.
%
\begin{align}
  \frac{\partial C_i}{\partial \theta_j}
  &= \frac{\delta_{ij}\,\exp(\theta_i)\,Z - \exp(\theta_i)\,\exp(\theta_j)}{Z^2}
  \nonumber\\
  &= \frac{\exp(\theta_i)}{Z}\Bigl(\delta_{ij} - \frac{\exp(\theta_j)}{Z}\Bigr)
  = C_i\,(\delta_{ij} - C_j).
  \label{eq:jacobian-derivation}
\end{align}

% SELF-CRITIQUE CHECKPOINT (step 1):
% 1. SIGN CHECK: Minus sign in (delta_ij - C_j) correct from quotient rule.
% 2. FACTOR CHECK: No spurious factors introduced.
% 3. CONVENTION CHECK: C_i = number fractions, sum to 1.
% 4. DIMENSION CHECK: All dimensionless. PASSED.

\subsection{Rank and null space}

The null space of $\JJ_\mathrm{sm}$ is spanned by $\ones = (1,\ldots,1)^T$:
%
\begin{equation}
  \JJ_\mathrm{sm}\,\ones
  = \mathrm{diag}(\CC)\,\ones - \CC\,\CC^T\,\ones
  = \CC - \CC\,\underbrace{(\CC^T\,\ones)}_{=\,1}
  = \CC - \CC = \mathbf{0}.
\end{equation}
%
$\therefore$ $\mathrm{rank}(\JJ_\mathrm{sm}) = D - 1$, corresponding to the
shift symmetry.

\subsection{Condition number analysis}

The nonzero eigenvalues of $\JJ_\mathrm{sm} = \mathrm{diag}(\CC) - \CC\CC^T$
are the eigenvalues of $\mathrm{diag}(\CC)$ projected onto the orthogonal
complement of $\CC$.

\paragraph{Uniform composition $C_i = 1/D$.}
$\JJ_\mathrm{sm} = (1/D)(I - \ones\ones^T/D)$.  The $D-1$ nonzero
eigenvalues are all $1/D$.
%
\begin{equation}
  \sigma_\mathrm{min} = \sigma_\mathrm{max} = \frac{1}{D},
  \qquad
  \kappa(\JJ_\mathrm{sm})\big|_\mathrm{uniform} = 1.
\end{equation}

\paragraph{Extreme composition ($C_1 \to 1$, $C_{i\ne 1} \to 0$).}
The largest nonzero singular value is $\sigma_\mathrm{max} \approx C_1(1-C_1) \approx \epsilon$
where $\epsilon = 1 - C_1 \ll 1$.  The smallest nonzero singular value is
$\sigma_\mathrm{min} \sim \min_i C_i$, which can be exponentially small.
The condition number grows as $\kappa \sim C_\mathrm{max}/C_\mathrm{min}$.

\subsection{Inverse mapping}

Given $\CC \in \Delta^{D-1}$ (with $C_i > 0$), the inverse is:
%
\begin{equation}
  \label{eq:softmax-inverse}
  \theta_i = \log(C_i) + c,
\end{equation}
%
where $c \in \mathbb{R}$ is an arbitrary constant (the shift freedom).
A common choice is $c = 0$, yielding $\theta_i = \log C_i$.  This
matches the \texttt{\_pack\_params} implementation in
\texttt{joint\_optimizer.py} (line 523).

\subsection{Hessian structure}

For an objective $L(\CC(\ttheta))$, the chain rule gives:
%
\begin{equation}
  \frac{\partial L}{\partial \theta_i}
  = \sum_j \frac{\partial L}{\partial C_j}\,
    C_j\,(\delta_{ij} - C_i)
  = C_i\,\Bigl(\frac{\partial L}{\partial C_i}
    - \sum_j C_j\,\frac{\partial L}{\partial C_j}\Bigr).
  \label{eq:softmax-gradient}
\end{equation}
%
% DIMENSION CHECK: [dL/dtheta] = [L] * [dim-less] = [L]. PASSED.
%
This is the ``centered gradient'': the raw gradient minus its
$\CC$-weighted mean.  JAX computes this automatically via autodiff.


%% ====================================================================
\section{ILR (Isometric Log-Ratio) Parameterization}
\label{sec:ilr}
%% ====================================================================

\subsection{Centered log-ratio (CLR)}

The CLR transform maps $\CC \in \Delta^{D-1}$ to the hyperplane
$\sum_i \mathrm{clr}_i = 0$ in $\mathbb{R}^D$:
%
\begin{equation}
  \label{eq:clr}
  \mathrm{clr}(\CC)_i = \log C_i - \frac{1}{D}\sum_{j=1}^D \log C_j.
\end{equation}
%
% DIMENSION CHECK: log(dim-less) = dim-less.  [clr_i] = dim-less.  PASSED.

\subsection{ILR via Helmert basis}

The ILR transform projects CLR coordinates onto a $(D{-}1)$-dimensional
orthonormal basis, removing the redundant zero-sum constraint:
%
\begin{equation}
  \label{eq:ilr}
  \zz = \mathrm{clr}(\CC)\,\VV \in \mathbb{R}^{D-1},
\end{equation}
%
where $\VV \in \mathbb{R}^{D \times (D-1)}$ is the Helmert contrast
matrix~\cite{egozcue-2003}.

\paragraph{Helmert basis construction.}
For $j = 1, \ldots, D-1$:
%
\begin{equation}
  \label{eq:helmert}
  V_{ij} = \begin{cases}
    \dfrac{1}{\sqrt{j(j+1)}}     & \text{if } i < j+1, \\[6pt]
    \dfrac{-j}{\sqrt{j(j+1)}}    & \text{if } i = j+1, \\[4pt]
    0                             & \text{if } i > j+1.
  \end{cases}
\end{equation}
%
(Using 1-based indexing for components, matching the codebase
\texttt{\_helmert\_basis} in \texttt{closure.py}.)

\paragraph{Orthonormality.}
$\VV^T\VV = I_{D-1}$ by construction.  Each column of $\VV$ sums to
zero: $\ones^T \VV = \mathbf{0}^T$.  This ensures the ILR coordinates are
orthogonal to the redundant CLR direction.

\paragraph{Verification for $D = 3$.}
%
\begin{equation}
  \VV = \begin{pmatrix}
    1/\sqrt{2} & 1/\sqrt{6} \\
    -1/\sqrt{2} & 1/\sqrt{6} \\
    0 & -2/\sqrt{6}
  \end{pmatrix}.
\end{equation}
%
Check: $\VV^T\VV = \begin{pmatrix} 1 & 0 \\ 0 & 1 \end{pmatrix}$. $\checkmark$

\subsection{Inverse ILR}

Given $\zz \in \mathbb{R}^{D-1}$, recover $\CC$:
%
\begin{align}
  \mathrm{clr} &= \zz\,\VV^T \in \mathbb{R}^D, \label{eq:ilr-inv-clr}\\
  C_i &= \frac{\exp(\mathrm{clr}_i)}{\sum_j \exp(\mathrm{clr}_j)}.
  \label{eq:ilr-inv}
\end{align}
%
Note: the inverse naturally uses the softmax of the CLR coordinates.
This matches \texttt{ilr\_inverse} in \texttt{closure.py} (line 123-124).

\subsection{Jacobian}

The full ILR Jacobian $\JJ_\mathrm{ilr} = \partial \CC / \partial \zz
\in \mathbb{R}^{D \times (D-1)}$ is the composition:
%
\begin{equation}
  \label{eq:ilr-jacobian}
  \boxed{
    \JJ_\mathrm{ilr}
    = \bigl(\mathrm{diag}(\CC) - \CC\,\CC^T\bigr)\,\VV
    = \JJ_\mathrm{sm}\,\VV.
  }
\end{equation}
%
% Derivation: dC_i/dz_k = sum_j (dC_i/d(clr_j)) * V_{jk}.
% The CLR-to-simplex map is the softmax of clr, so dC_i/d(clr_j) = J^sm_{ij}.
% Therefore dC_i/dz_k = sum_j J^sm_{ij} V_{jk} = (J_sm * V)_{ik}.

% DIMENSION CHECK: [J_ilr] = dim-less * dim-less = dim-less. PASSED.

\paragraph{Rank.}
$\JJ_\mathrm{ilr}$ has rank $D-1$ because $\JJ_\mathrm{sm}$ has rank
$D-1$ and $\VV$ has $D-1$ columns with $\ones^T\VV = \mathbf{0}^T$
(so $\VV$ maps entirely into the row space of $\JJ_\mathrm{sm}$).

\subsection{Condition number analysis}

\paragraph{Uniform composition.}
$\JJ_\mathrm{sm} = (1/D)(I - \ones\ones^T/D)$ and $\VV^T\VV = I_{D-1}$,
so the singular values of $\JJ_\mathrm{ilr} = (1/D)(I - \ones\ones^T/D)\VV$
are all $1/D$ (the Helmert rotation preserves the spectrum on the
orthogonal complement of $\ones$).
%
\begin{equation}
  \kappa(\JJ_\mathrm{ilr})\big|_\mathrm{uniform} = 1.
\end{equation}

\paragraph{Near boundaries ($C_i \to 0$).}
The ILR transform requires $C_i > 0$ strictly, since $\log(0) = -\infty$.
As $C_i \to 0$, the CLR coordinates diverge logarithmically, and the
Jacobian develops a singularity.  This is \emph{worse} than softmax:
softmax naturally excludes $C_i = 0$ (since $\exp(\theta_i) > 0$ always),
and the approach is smooth as $\theta_i \to -\infty$.  ILR requires
explicit clamping ($C_i \ge \epsilon$) to avoid numerical breakdown.

% SELF-CRITIQUE CHECKPOINT (step 2):
% 1. SIGN CHECK: Minus in J_sm = diag(C) - CC^T verified in Sec 2.
% 2. FACTOR CHECK: No spurious factors. Helmert normalization 1/sqrt(j(j+1)) correct.
% 3. CONVENTION CHECK: Compositions are number fractions (sum=1). Consistent.
% 4. DIMENSION CHECK: All quantities dimensionless. PASSED.


%% ====================================================================
\section{Comparison}
\label{sec:comparison}
%% ====================================================================

\begin{table}[h]
\centering
\caption{Softmax vs.\ ILR for CF-LIBS compositional closure.}
\label{tab:comparison}
\begin{tabular}{@{}lll@{}}
\toprule
\textbf{Property}        & \textbf{Softmax}                    & \textbf{ILR} \\
\midrule
Parameters               & $D$ (overparameterized)              & $D-1$ (minimal) \\
Domain                   & $\mathbb{R}^D \to \Delta^{D-1}$     & $\mathbb{R}^{D-1} \to \Delta^{D-1}$ \\
Boundary behavior        & Smooth ($C_i \to 0$ as $\theta_i \to -\infty$) & Singular ($\log$ diverges at $C_i = 0$) \\
Jacobian rank            & $D-1$ (null space: $\ones$)          & $D-1$ (full rank) \\
$\kappa$ at uniform      & $1$                                  & $1$ (isometric) \\
$\kappa$ at boundary     & $\sim C_\mathrm{max}/C_\mathrm{min}$ & $\sim 0$ ($\log$ singularity) \\
JIT compatible           & Yes (\texttt{exp}, \texttt{sum})     & Yes (\texttt{log}, \texttt{matmul}, \texttt{exp}) \\
Autodiff                 & Clean (no branches)                  & Clean (no branches) \\
Implementation           & 1 line (\texttt{jax.nn.softmax})     & $\sim$10 lines (Helmert + CLR + project) \\
Codebase                 & \texttt{joint\_optimizer.py}         & \texttt{closure.py} \\
\bottomrule
\end{tabular}
\end{table}


%% ====================================================================
\section{Gradient Flow Analysis}
\label{sec:gradient}
%% ====================================================================

For an objective function $L(\CC)$ optimized via L-BFGS-B over
$\ttheta$ (softmax) or $\zz$ (ILR):

\paragraph{Softmax gradient.}
From Eq.~\eqref{eq:softmax-gradient}:
%
\begin{equation}
  \frac{\partial L}{\partial \theta_i}
  = C_i\,\biggl(\frac{\partial L}{\partial C_i}
    - \sum_j C_j\,\frac{\partial L}{\partial C_j}\biggr).
\end{equation}
%
The gradient magnitude scales as
$\|\partial L/\partial\ttheta\| \sim \max(C_i)\,\|\partial L/\partial\CC\|$,
since the spectral norm of $\JJ_\mathrm{sm}$ is $\max(C_i)(1 - \max(C_i))
\le \max(C_i)$.

\paragraph{ILR gradient.}
%
\begin{equation}
  \frac{\partial L}{\partial \zz}
  = \VV^T\,\bigl(\mathrm{diag}(\CC) - \CC\CC^T\bigr)\,
    \frac{\partial L}{\partial \CC}.
  \label{eq:ilr-gradient}
\end{equation}
%
This is the same Jacobian as softmax, composed with the Helmert rotation.
The scaling is identical.

\paragraph{Practical conclusion.}
Gradient flow is comparable for both parameterizations.  For typical
LIBS compositions (major element at 50--90\%, minor elements 0.1--10\%),
the Jacobian condition number $\kappa \sim C_\mathrm{max}/C_\mathrm{min}
\sim 10$--$1000$.  Softmax has slightly better conditioning near simplex
boundaries because it avoids the logarithmic singularity.


%% ====================================================================
\section{Limiting Cases}
\label{sec:limits}
%% ====================================================================

\subsection{$D = 1$: Trivial (single element)}

\textbf{Softmax:}
$C_1 = \exp(\theta_1)/\exp(\theta_1) = 1$, identically, for any
$\theta_1$.  The Jacobian is $J_{11} = C_1(1 - C_1) = 0$
(zero-dimensional simplex). $\checkmark$

\textbf{ILR:}
$\zz$ is the empty vector ($\mathbb{R}^0$).  The Helmert basis $\VV$
is empty.  $C_1 = 1$.  $\checkmark$

\subsection{$D = 2$: Sigmoid}

\textbf{Softmax:}
%
\begin{align}
  C_1 &= \frac{\exp(\theta_1)}{\exp(\theta_1) + \exp(\theta_2)}
       = \frac{1}{1 + \exp(\theta_2 - \theta_1)}
       = \sigma(\theta_1 - \theta_2),
  \label{eq:sigmoid}
\end{align}
%
where $\sigma(x) = 1/(1+e^{-x})$ is the logistic sigmoid.
$C_2 = 1 - C_1 = \sigma(\theta_2 - \theta_1)$. $\checkmark$

\textbf{ILR:}
$\zz = z_1 \in \mathbb{R}^1$.  The Helmert basis is
$\VV = (1/\sqrt{2},\;-1/\sqrt{2})^T$.
%
\begin{align}
  \mathrm{clr}_1 &= z_1/\sqrt{2}, \quad
  \mathrm{clr}_2 = -z_1/\sqrt{2}. \nonumber\\
  C_1 &= \frac{\exp(z_1/\sqrt{2})}{\exp(z_1/\sqrt{2}) + \exp(-z_1/\sqrt{2})}
       = \sigma(\sqrt{2}\,z_1).
\end{align}
%
Setting $\theta_1 - \theta_2 = \sqrt{2}\,z_1$ shows the two
parameterizations are related by a simple scaling. $\checkmark$

\subsection{Uniform: $\theta_i = c$ for all $i$}

\textbf{Softmax:}
$C_i = e^c / (D\,e^c) = 1/D$ for all $i$. $\checkmark$

\textbf{ILR:}
$\zz = \mathbf{0}$ maps to $\mathrm{clr} = \mathbf{0}$, hence
$C_i = 1/D$. $\checkmark$

\subsection{ILR round-trip identity}

For any $\CC \in \Delta^{D-1}$ with $C_i > 0$:
%
\begin{equation}
  \mathrm{ilr\_inverse}\bigl(\mathrm{ilr}(\CC)\bigr) = \CC.
\end{equation}
%
\textbf{Proof:}
$\mathrm{ilr}(\CC) = \mathrm{clr}(\CC)\,\VV$.
Then $\mathrm{ilr\_inverse}(\zz) = \mathrm{softmax}(\zz\,\VV^T)
= \mathrm{softmax}(\mathrm{clr}(\CC)\,\VV\VV^T)$.
Since $\VV\VV^T$ is the projector onto the orthogonal complement of
$\ones/\sqrt{D}$ in $\mathbb{R}^D$, and $\mathrm{clr}(\CC)$ already lies
in this complement ($\ones^T\,\mathrm{clr} = 0$), we have
$\mathrm{clr}(\CC)\,\VV\VV^T = \mathrm{clr}(\CC)$.
Therefore
$\mathrm{softmax}(\mathrm{clr}(\CC)) = \mathrm{softmax}(\log\CC - \bar{\log\CC}\,\ones)
= \mathrm{softmax}(\log\CC) = \CC$.
$\checkmark$

This matches the \texttt{ilr\_inverse(ilr\_transform(simplex))} round-trip
in \texttt{closure.py} (line 452--453).

% SELF-CRITIQUE CHECKPOINT (step 3):
% 1. SIGN CHECK: All signs verified in sigmoid derivation and ILR D=2 case.
% 2. FACTOR CHECK: sqrt(2) factor in ILR D=2 case verified from Helmert basis.
% 3. CONVENTION CHECK: Compositions are number fractions (sum=1). Consistent.
% 4. DIMENSION CHECK: All dimensionless throughout. PASSED.


%% ====================================================================
\section{Recommendation and Validation}
\label{sec:recommendation}
%% ====================================================================

\subsection{Primary recommendation}

\textbf{Softmax for the GPU pipeline} because:
\begin{enumerate}
  \item No boundary singularity: graceful behavior as $C_i \to 0$.
  \item Simpler implementation: \texttt{jax.nn.softmax} (1 line).
  \item Already implemented and tested in \texttt{joint\_optimizer.py}.
  \item Well-conditioned for typical LIBS compositions ($D = 2$--$15$
    elements, major element at 50--90\%).
\end{enumerate}

\textbf{ILR for diagnostic/analysis} where Aitchison geometry matters:
\begin{itemize}
  \item Visualizing composition trajectories (ILR preserves Euclidean
    distances in compositional space).
  \item PCA in composition space (standard PCA in ILR coordinates is
    equivalent to compositional PCA).
  \item Statistical analysis of compositional uncertainty distributions.
\end{itemize}

\subsection{Validation criteria (Phase 4)}

\begin{enumerate}
  \item $\sum_i \mathrm{softmax}(\ttheta)_i = 1$ to machine precision
    ($< 10^{-15}$) for random $\ttheta \in [-100, 100]^D$.
  \item $\mathrm{ilr\_inverse}(\mathrm{ilr}(\CC)) = \CC$ to machine
    precision for random $\CC$ on the simplex.
  \item The Jacobian $\JJ_\mathrm{sm}$ has exactly $D-1$ nonzero
    singular values (verified numerically for $D = 2, \ldots, 15$).
  \item JAX \texttt{grad(loss)(theta)} matches finite-difference
    gradient to relative error $< 10^{-6}$.
\end{enumerate}


%% ====================================================================
\begin{thebibliography}{9}

\bibitem{egozcue-2003}
J.J.~Egozcue, V.~Pawlowsky-Glahn, G.~Mateu-Figueras, and
C.~Barcel\'o-Vidal,
``Isometric logratio transformations for compositional data analysis,''
\emph{Math. Geology} \textbf{35}(3), 279--300 (2003).

\bibitem{aitchison-1986}
J.~Aitchison,
\emph{The Statistical Analysis of Compositional Data},
Chapman and Hall, London (1986).

\end{thebibliography}

\end{document}
